\thispagestyle{plain}
\typeout{HVRPAGE|1|\thepage}
\pdfbookmark[0]{Title and abstract}{title}
\begin{center}
{\LARGE\bfseries Hilbert--Vector--Resolvent Theory\par}
\vspace{8pt}
{\large Positive Moments and Two Conditional Routes\par
to the Riemann Hypothesis\par}
\vspace{10pt}
Preprint \textperiodcentered\ Expanded draft 0.1\par
28 October 2026
\end{center}
\vspace{5pt}
\begin{abstract}
We develop an introductory framework connecting scalar coefficients,
distinguished vectors in Hilbert space, positive moment measures, and
analytic resolvents. For a positive contraction $T$ and a unit vector
$v$, the moments $\langle v,T^nv\rangle$ have nonnegative Bernstein
weights; conversely, normalized Bernstein positivity produces a
probability measure on $[0,1]$ and a concrete multiplication-operator
model. The associated scaled resolvent is holomorphic on a slit plane
and has nonpositive imaginary part in the upper half-plane. We give
the finite algebra, moment-limit construction, and analytic
continuation argument underlying this correspondence.

The arithmetic application concerns the actual normalized moments
of the centered Riemann $\xi$-function. Their universal zero-cancellation
condition, denoted $H00$, yields two conditional routes to the literal
Riemann hypothesis: one through spectral growth of Bernstein variation,
and one through a positive measure and its resolvent. The existing Lean
development proves both implications and the corresponding equivalences;
the present text supplies a broader mathematical exposition and a source
concordance. The concrete Hilbert-space packaging is presented as
classical mathematics, separately from the declarations already audited.
The hypothesis $H00$ itself remains unproved. Neither an unconditional
proof of the Riemann hypothesis nor a finite accepted certificate is
asserted.
\end{abstract}

\noindent\textbf{Keywords.} Hilbert space; positive contraction;
Hausdorff moments; Bernstein weights; scalar resolvent; Riemann
hypothesis; formal verification.

\medskip
\noindent\textbf{Scope of this draft.}
``Hilbert--Vector--Resolvent theory'' is the organizing name used here
for the project. The elementary operator and moment correspondences
belong to classical analysis. The purpose is to state the project's
objects and proof dependencies precisely, to explain the two existing
conditional constructions, and to isolate the exact remaining
identification and admissibility problems. No priority claim for the
classical correspondences is intended.

\medskip
\noindent\textbf{Source baseline.}
The mathematical source and verification record are available in the
\href{https://github.com/randogelist/Hilber-Vector-Resolvent-Theory}{project repository}.
The frozen project anchor is \code{ZETA_HYPERBOLIC_ANCHOR_5}.

\PaperPage{2}{Contents and reading guide}
\begin{tabular}{@{}p{0.82\linewidth}r@{}}
\toprule
\textbf{Part} & \textbf{Pages}\\
\midrule
Abstract, notation, principal claims, and introduction & 1--6\\[5pt]
1. Hilbert-vector models & 7--14\\
\quad Coefficients, defects, positivity, cyclicity, examples & \\
2. Bernstein positivity and the moment measure & 15--22\\
\quad Finite identities, shifted moments, positive functionals & \\
3. Scalar resolvents and analytic zero exclusion & 23--30\\
\quad Slit domain, derivatives, signs, continuation, poles & \\
4. Actual $\xi$-moments and the two RH routes & 31--38\\
\quad Normalization, spectral growth, common converse & \\
5. Weighted models and the unfinished frontier & 39--44\\
\quad Kernel models, exact identification, finite certificates & \\
Appendix A. Verification evidence and exact endpoints & 45--46\\
Appendix B. Finite prefixes and weak convergence & 47--48\\
References & 49\\
Appendix C. Source concordance and reproducibility & 50\\
\bottomrule
\end{tabular}

\medskip
The mathematical narrative runs from general objects to the actual
arithmetic application. Part~1 fixes the meaning of a Hilbert-vector
model and explains why one vector sees a cyclic subspace rather than
an arbitrary ambient operator. Part~2 works directly with real
sequences. It derives positivity from finite identities before passing
to an infinite moment measure. Part~3 studies the analytic function
determined by that measure, with all sign conventions fixed explicitly.

Part~4 specializes the construction to the actual centered
$\xi$-function. Its retained spectral infrastructure is stated with its
role and provenance; the shorter resolvent route is then assembled
from the preceding arguments. Part~5 returns to the broader program
and states the missing equality needed to connect the weighted model
to the archived Widder family. The appendices separate compiler
evidence from mathematical interpretation and record useful limits of
finite testing.

Readers interested primarily in the conditional RH result can begin
with the statement on page~4 and Part~4, consulting Parts~2 and~3 for
the second route. Readers interested in a rigorous general introduction
can read Parts~1--3 independently of the arithmetic application.
The bibliography supplies classical background; the source concordance
gives the more granular project-level correspondence.

Each numbered subsection is a self-contained stage of the argument.
Proofs use standard Hilbert-space, measure-theoretic, and complex-analytic
results where named. When a substantial theorem is imported from the
retained formal development, it is identified as such rather than
represented as an elementary consequence of the surrounding exposition.

\PaperPage{3}{Notation and conventions}
\begin{tabular}{@{}p{0.21\linewidth}p{0.73\linewidth}@{}}
\toprule
Symbol & Meaning\\
\midrule
$\mathcal H,T,v$ & Complex Hilbert space, bounded operator, distinguished vector.\\[4pt]
$\langle x,y\rangle$ & Inner product conjugate-linear in $x$, linear in $y$.\\[4pt]
$u_n$ & Real normalized moment; $u_0=1$ when positivity is used.\\[4pt]
$w_{N,j}$ & Bernstein weight of degree $N$, for $0\le j\le N$.\\[4pt]
$V_N,C_N$ & Row variation and row cancellation, respectively.\\[4pt]
$\sigma$ & Borel probability measure on $[0,1]$.\\[4pt]
$\kappa$ & Positive scaling parameter; actual trace moment in Part~4.\\[4pt]
$S_\kappa$ & $\kappa\int_{[0,1]}(1+\kappa qz)^{-1}\,d\sigma(q)$.\\[4pt]
$\mathbb H$ & Upper half-plane $\{z\in\mathbb C:\operatorname{Im}z>0\}$.\\[4pt]
$\xi,\Theta$ & Entire Riemann $\xi$ and its even centered descent.\\[4pt]
$R$ & The germ $\Theta'/\Theta$ at the origin.\\[4pt]
$H00$ & Semantic admissibility of the actual zero-budget candidate.\\
\bottomrule
\end{tabular}

\medskip
All sequence indices start at zero. The binomial coefficient is the
ordinary integer binomial coefficient. Finite sums over an empty set
are zero. The grid point at degree zero is set to $q_{0,0}=0$; for
$N>0$, the grid is $q_{N,j}=j/N$. This convention avoids interpreting
$0/0$ when defining Bernstein averages.

For a bounded self-adjoint operator, $0\le T\le I$ means that
$\langle x,Tx\rangle$ and $\langle x,(I-T)x\rangle$ are nonnegative
for every $x$. The term \emph{positive contraction} includes both
self-adjointness and these inequalities. A bare contraction need not
have real or positive scalar coefficients.

The two scalar resolvent conventions must be kept distinct:
\[
 F(z)=\langle v,(I-zT)^{-1}v\rangle,\qquad
 S_\kappa(z)=\kappa\langle v,(I+\kappa zT)^{-1}v\rangle.
\]
For a positive model, $F$ has nonnegative imaginary part on
$\mathbb H$, while $S_\kappa$ has nonpositive imaginary part there.
The inequalities need not be strict; the measure $\delta_0$ gives a
constant resolvent. ``Anti-Herglotz'' will mean the latter nonpositive
sign condition, with the convention always displayed when used.

An equality of germs means equality throughout some neighborhood of
the stated point. Holomorphy on a domain does not assert that this
domain is maximal. In particular, the universal slit used for
$S_\kappa$ is a common domain guaranteed for every probability measure
on $[0,1]$, not the singular set of every particular example.

\PaperPage{4}{Principal claims and their logical scope}
The analytic framework has a general equivalence independent of the
zeta function. It is a classical form of the compact moment problem.
\begin{theorem}[Positive-model correspondence]\label{intro:correspondence}
Let $(u_n)_{n\ge0}$ be real and $u_0=1$. The following are equivalent:
\begin{enumerate}[label=(\roman*)]
\item Every Bernstein weight $w_{N,j}$ defined on page~15 is nonnegative.
\item There is a Borel probability measure $\sigma$ on $[0,1]$ such that
$u_n=\int q^n\,d\sigma(q)$ for all $n$.
\item There is a positive contraction $T$ on a complex Hilbert space
and a unit vector $v$ with $u_n=\langle v,T^nv\rangle$.
\end{enumerate}
The measure in (ii) is unique. One may choose the cyclic model
$\mathcal H=L^2(\sigma)$, $T=M_q$, $v=1$.
\end{theorem}
The implication (i)$\Rightarrow$(ii) is developed in Part~2;
(ii)$\Rightarrow$(iii) is the multiplication construction in Part~1;
and (iii)$\Rightarrow$(i) follows from the positive operators
$T^j(I-T)^{N-j}$. Uniqueness follows from polynomial density.

\begin{theorem}[Retained actual conditional endpoints]\label{intro:actual}
For the actual normalized $\xi$-moment sequence defined in Part~4,
the semantic condition $H00$ is equivalent to the literal Riemann
hypothesis. The retained development provides two forward
constructions of $H00\Rightarrow\mathrm{RH}$ and a common converse.
\end{theorem}
Here $H00$ is a universal assertion about all Bernstein rows. Its
finite candidate data are $(0,0)$, but its admissibility is not a
finite computation. The first route uses the spectral growth theorem
for row variation. The second constructs the measure, identifies the
actual logarithmic-derivative germ, and excludes nonreal zeros.
Precise endpoint types and their audit record appear in Appendix~A.

There are three distinct levels of claim in this manuscript.
The finite identities and actual conditional endpoints have source
counterparts covered by the existing audit. The full complex Hilbert
realization and some general explanatory lemmas are proved here as
ordinary mathematics, without a claim of new Lean coverage. Finally,
unconditional $H00$, checker soundness and completeness, and the
archived Widder identification remain goals. None is promoted to a
theorem by the existence of a model with similarly named coefficients.

This separation is mathematically consequential. A positive measure
for an arbitrary sequence gives a resolvent for that sequence. To
derive a statement about $\xi$, the sequence must be the actual Taylor
data of its logarithmic derivative. That identification is the link
which makes the second route an arithmetic argument.

\PaperPage{5}{Introduction: from scalar data to an operator model}
A sequence often enters analysis in several apparently different
forms: as the derivatives of a germ, as moments of a measure, or as
matrix coefficients of an operator orbit. The useful question is not
whether one can assign these descriptions formally, but which
identities and inequalities make them describe the same object.
The framework of this preprint begins with the chain
\[
 u_n=\int_0^1 q^n\,d\sigma(q)
     =\langle 1,M_q^n1\rangle_{L^2(\sigma)}.
\]
The moment measure supplies positivity and compact support. The
operator supplies norm estimates and geometric defects. The resolvent
packages all moments into a single analytic function. Each description
therefore exposes a different part of the same structure.

The finite entry point is the Bernstein basis
$\binom Nj q^j(1-q)^{N-j}$. If a normalized sequence acts positively
on every member of this basis, its row weights form probability
distributions on increasingly fine grids. The finite coarsening
identity measures how total variation changes from one row to the
next. Vanishing cancellation is precisely the condition that all
these signed rows are already positive. The passage to a measure then
requires an actual limit argument; row positivity alone is not a
license to interchange arbitrary infinite operations.

Once the measure exists, its compact support controls denominators:
\[
 S_\kappa(z)=\kappa\int_0^1\frac{d\sigma(q)}{1+\kappa qz},
 \qquad z\notin(-\infty,-1/\kappa].
\]
The imaginary part of the integrand has a fixed sign in the upper
half-plane. This sign becomes a zero-exclusion mechanism only after
the resolvent is identified with the logarithmic derivative of an
entire function. At a zero of positive multiplicity, that logarithmic
derivative has a pole with a positive integer residue. A local
downward approach to the pole contradicts the resolvent sign.

The actual $\xi$ application is an instance of this mechanism with
carefully normalized Taylor data. It also admits a different route
through exponential growth of Bernstein variation. The coexistence
of the two routes is useful: one makes the geometry of the canonical
zeros visible, while the other makes the positive measure and analytic
sign visible. Their common unresolved premise is the same exact
universal admissibility statement.

The introduction therefore aims at a disciplined synthesis. Classical
Hilbert and moment theory provide the language; the repository gives
specific actual constructions and audited endpoints; the remaining
program is expressed as explicit assertions whose proofs can be
sought independently.

\PaperPage{6}{Representation, identification, and the proof architecture}
Three operations organize the discussion. \emph{Representation}
constructs an object realizing given data. \emph{Identification}
proves that these data equal those of a prescribed analytic function.
\emph{Closure} converts the resulting analytic information into a
statement about zeros. Confusing any two of these operations can hide
a missing hypothesis.

For representation, Parts~1 and~2 connect four descriptions:
\[
 \begin{gathered}
 \text{nonnegative Bernstein rows}
 \ \Longleftrightarrow\ \text{positive moment measure on }[0,1],\\
 \text{positive moment measure}
 \ \Longleftrightarrow\ \text{cyclic positive-contraction model}.
 \end{gathered}
\]
The second equivalence is understood up to unitary equivalence of
cyclic models. It is not a uniqueness assertion for all possible
ambient spaces containing an unobserved orthogonal summand.
The resulting resolvent has prescribed Taylor coefficients and a
half-plane sign, independently of any arithmetic interpretation.

For identification, let $\Theta(0)\ne0$ and assume that the actual
derivatives satisfy
\[
 \left(\frac{\Theta'}{\Theta}\right)^{(n)}(0)
   =(-1)^n n!\kappa^{n+1}u_n.
\]
The same formula holds for $S_\kappa^{(n)}(0)$. Analyticity makes
the two germs equal. Their equality is propagated using the holomorphic
defect $\Theta'-S_\kappa\Theta$ on the upper half-plane; dividing by
$\Theta$ before excluding its zeros would require unnecessary care.
For closure, the positive-multiplicity pole of $\Theta'/\Theta$ is
incompatible with the anti-Herglotz sign. Conjugation symmetry then
handles the lower half-plane.

The first arithmetic route uses a different intermediate invariant:
\[
 H00\Longrightarrow V_N=1\ (\forall N)
 \Longrightarrow \text{subexponential variation}
 \Longrightarrow\mathrm{RH}.
\]
The last implication uses the retained actual spectral extraction
theorem. The second route replaces that forward spectral-growth step
by the positive-measure construction and analytic closure. It shares
normalization, direct positivity lemmas, and the converse from RH.

The classical background is the compact Hausdorff moment problem,
the spectral theory of bounded self-adjoint operators, and elementary
complex analysis \cite{hausdorff,conway,rudin}. The arithmetic
normalization is compatible with the classical $\xi$-function
\cite{titchmarsh}. Formal proof infrastructure is supplied by Lean
and mathlib \cite{lean4,mathlib}; the project's precise theorem scope
is fixed by its source and audit records \cite{source}.
